Bij SVM met Shin en eigen categorieën geeft het husselen van de buiging van de duim de grootste gemiddelde F1-daling (0,08 procentpunt). Bij LightGBM is dat de spreiding tussen middel- en ringvinger (0,03 procentpunt). Dit laat zien welke informatie deze modellen gebruiken. Een kleine daling bewijst niet dat een kenmerk nutteloos is: andere kenmerken kunnen dezelfde informatie leveren. De detailresultaten staan in bijlage B.
